\renewcommand\qedsymbol{Q.E.F}

\begin{con}
To draw a straight line perpendicular to a given infinite straight line from a given point not on it.
\end{con}

Let $\overline{AB}$ be the given infinite straight line, and $C$ the given point not on it. It is required to draw a straight line perpendicular to the given infinite straight line $\overline{AB}$ from the given point $C$.

\begin{figure}[H]
\centering
	\begin{tikzpicture}
		\tkzDefPoint(-1,0){A}    		\tkzLabelPoint[left](A){$A$}	
		\tkzDefPoint(5,0){B}    		\tkzLabelPoint[right](B){$B$}
		\tkzDefPoint(2,2){C}			\tkzLabelPoint[above](C){$C$}
		\tkzDrawSegment(A,B)
		\tkzDrawPoints[fill=gray](A,B,C)
	\end{tikzpicture}
	\caption{}
\end{figure}

\begin{enumerate}

	\item[]

	\item Take an arbitrary point $D$ on the other side of the straight line $\overline{AB}$.

\begin{figure}[H]
\centering
	\begin{tikzpicture}
		\tkzDefPoint(-1,0){A}    		\tkzLabelPoint[left](A){$A$}	
		\tkzDefPoint(5,0){B}    		\tkzLabelPoint[right](B){$B$}
		\tkzDefPoint(2,2){C}			\tkzLabelPoint[above](C){$C$}
		\tkzDefPoint(2.5,-0.8){D}	\tkzLabelPoint[below](D){$D$}
		\tkzDrawSegment(A,B)
		\tkzDrawPoints[fill=gray](A,B,C,D)
	\end{tikzpicture}
	\caption{}
\end{figure}

\clearpage

	\item Describe $\mathscr{C}(C;CD)$.

\begin{figure}[H]
\centering
	\begin{tikzpicture}
		\tkzDefPoint(-1,0){A}    		\tkzLabelPoint[left](A){$A$}	
		\tkzDefPoint(5,0){B}    		\tkzLabelPoint[right](B){$B$}
		\tkzDefPoint(2,2){C}			\tkzLabelPoint[above](C){$C$}
		\tkzDefPoint(2.5,-0.8){D}		
		\tkzInterLC(A,B)(C,D)	\tkzGetPoints{E}{G}
		\tkzLabelPoints[below](D,E,G)
		\tkzDrawCircle(C,D)
		\tkzDrawSegments(A,B)
		\tkzDrawPoints[fill=gray](A,B,C,D,E,G)
	\end{tikzpicture}
	\caption{}
\end{figure}

	\item Bisect the straight line $\overline{EG}$ at $H$, and join the straight lines $\overline{CG}$, $\overline{CH}$, and $\overline{CE}$.\hfill\textcolor{red}{I.10}

\begin{figure}[H]
\centering
	\begin{tikzpicture}
		\tkzDefPoint(-1,0){A}    		\tkzLabelPoint[left](A){$A$}	
		\tkzDefPoint(5,0){B}    		\tkzLabelPoint[right](B){$B$}
		\tkzDefPoint(2,2){C}			\tkzLabelPoint[above](C){$C$}
		\tkzDefPoint(2.5,-0.8){D}		
		\tkzInterLC(A,B)(C,D)	\tkzGetPoints{E}{G}
		\tkzDefMidPoint(E,G)	\tkzGetPoint{H}
		\tkzLabelPoints[below](D,E,G,H)
		\tkzCompass[blue,thick](C,E)
		\tkzCompass[blue, thick](C,G)
		\tkzDrawSegments(A,B C,G C,H C,E)
		\tkzDrawPoints[fill=gray](A,B,C,D,E,G,H)
		\tkzMarkRightAngle[red](E,H,C)
	\end{tikzpicture}
	\caption{}
\end{figure}

\end{enumerate}

\begin{proof}

\begin{itemize}

\item[]

\item Since 
\[\overline{GH} = \overline{HE}\]

\item and $\overline{HC}$ is common, $\therefore$ the two sides $\overline{GH}$ and $\overline{HC}$ equal the two sides $\overline{EH}$ and $\overline{HC}$ respectively

\clearpage

\item and the bases 
\[\overline{CG} = \overline{CE}\] 

\item[$\therefore$] 
\[\angle{CHG} = \angle{EHC}\] 
and they are adjacent angles.

\item When a straight line standing on a straight line makes the adjacent angles equal to one another, each of the equal angles is right, and the straight line standing on the other is called a perpendicular to that on which it stands.\hfill\textcolor{red}{ I.Def.10}

\item[$\therefore$] $\overline{CH}$ has been drawn perpendicular to the given infinite straight line $\overline{AB}$ from the given point $C$ not on it.

\end{itemize}

\end{proof}

\clearpage
